% GCC compilation pipeline, marking the point where this work's synthesis
% happens (the RTL expansion pass, where `define_expand` bodies run).
%
% TikZ port of figures/pipeline.typ.  Conventions:
%   - solid boxes / arrows : the representations and the passes between them
%   - dashed grey boxes    : machine-description constructs acting at that step
\begin{figure}[htb]
\centering
\begin{SingleSpace}
\begin{tikzpicture}[
  x=0.74cm, y=0.74cm,
  font=\sffamily\fontsize{7}{8}\selectfont,
  stage/.style={draw, fill=black!4, inner sep=0pt, align=center,
                line width=0.5pt},
  flow/.style={-{Straight Barb[scale=0.55]}, line width=0.5pt},
  note/.style={draw=black!45, dashed, line width=0.4pt},
  notetext/.style={anchor=north west, text=black!55, align=left,
                   font=\sffamily\fontsize{5.5}{6.6}\selectfont},
  passlbl/.style={anchor=east, align=right,
                  font=\sffamily\fontsize{5.5}{6.6}\selectfont},
  bracket/.style={draw=black!50, line width=0.5pt},
]

\def\bxa{2.2}
\def\bxb{6.4}
\def\ax{4.3}

\newcommand{\Stage}[2]{%
  \path[stage] (\bxa,{#1-0.45}) rectangle (\bxb,{#1+0.45});
  \node at ({(\bxa+\bxb)/2},#1) {#2};}
% Vertical arrow between two stages, with the pass name to its left.
\newcommand{\Step}[3]{%
  \draw[flow] (\ax,{#1-0.45}) -- (\ax,{#2+0.45});
  \node[passlbl] at ({\ax-0.25},{(#1+#2)/2}) {#3};}
% Dashed note box on the right, with a pointer into the flow.
\newcommand{\Note}[4]{%
  \path[note] (10.2,#1) rectangle (18.1,#2);
  \node[notetext] at (10.4,{#1-0.2}) {#4};
  \draw[note, -{Straight Barb[scale=0.5]}] (10.2,#3) -- ({\ax+0.15},#3);}
% Phase bracket down the left-hand side.
\newcommand{\Phase}[3]{%
  \draw[bracket] (-0.5,#1) -- (-0.5,#2);
  \draw[bracket] (-0.5,#1) -- (-0.3,#1);
  \draw[bracket] (-0.5,#2) -- (-0.3,#2);
  \node[rotate=-90, text=black!60,
        font=\sffamily\fontsize{5.5}{6.6}\selectfont]
    at (-0.85,{(#1+#2)/2}) {#3};}

% --- the representations -------------------------------------------------
\Stage{0}{C source}
\Stage{-1.9}{GENERIC (AST)}
\Stage{-3.9}{GIMPLE, SSA form}
\Stage{-6.3}{RTL, pseudo-registers}
\Stage{-8.4}{RTL, hard registers}
\Stage{-10.2}{assembly (\texttt{.s})}
\Stage{-12.0}{ELF object}

% --- the passes between them ---------------------------------------------
\Step{0}{-1.9}{parse}
\Step{-1.9}{-3.9}{gimplify,\\SSA construction}
\Step{-3.9}{-6.3}{\textbf{RTL expansion}}
\Step{-6.3}{-8.4}{combine, register\\allocation (IRA/LRA),\\post-reload splits}
\Step{-8.4}{-10.2}{final: output\\templates}
\Step{-10.2}{-12.0}{\texttt{as} (binutils)}

\node[anchor=west, text=black!60, align=left,
      font=\sffamily\fontsize{5.5}{6.6}\selectfont\itshape]
  at ({\bxb+0.3},-3.9) {target-independent\\optimisation passes};

% --- phase brackets -------------------------------------------------------
\Phase{0.5}{-4.4}{front end + middle end}
\Phase{-4.9}{-10.7}{back end (target dependent)}

% --- what the machine description contributes, and where ------------------
\path[note] (10.2,0.5) rectangle (18.1,-1.15);
\node[notetext] at (10.4,0.3) {%
  \texttt{rvscN.h} \texttt{CC1\_SPEC} injects \texttt{-mno-shift},
  \texttt{-mno-xor}, \dots\\
  which clear the \texttt{riscv.opt} flags \texttt{TARGET\_SHIFT},\\
  \texttt{TARGET\_XOR}, \dots\ tested by every condition below};
\draw[note, -{Straight Barb[scale=0.5]}] (13.4,-1.15) -- (13.4,-3.7);

\Note{-3.7}{-5.7}{-5.1}{%
  \texttt{define\_expand} --- \textbf{where the synthesis happens.}\\
  With \texttt{-mno-shift}, \texttt{ashlsi3} expands to a chain of
  \texttt{add}\\
  instructions and calls \texttt{DONE}, so the \texttt{sll} pattern is never\\
  reached: the operation ceases to exist as RTL.}
\Note{-6.5}{-8.0}{-7.35}{%
  \texttt{define\_insn\_and\_split} --- splits that must run\\
  after register allocation, such as replacing\\
  \texttt{andi rd,rs,0xff} by \texttt{li} + \texttt{and}.}
\Note{-8.5}{-10.0}{-9.3}{%
  \texttt{define\_insn} --- output templates. A pattern whose\\
  condition is false (\texttt{TARGET\_SHIFT} on rvsc1) is\\
  invisible to instruction selection.}

\end{tikzpicture}

\vspace{0.4em}
{\footnotesize\color{black!60}\raggedright
The RTL expansion pass is where the synthesis of this work happens: a
\texttt{define\_expand} guarded by a \texttt{-mno-} flag emits a replacement
sequence and calls \texttt{DONE}, so the absent instruction never enters RTL and
no later pass can reintroduce it.\par}

\end{SingleSpace}
\caption{GCC compilation pipeline and the machine-description construct acting
         at each step}
\label{fig:pipeline}
\end{figure}
